\begingroup\color{blue}
\subsection{Sensitivity to Measurement Noise and External Forces}
Fig.~\ref{fig:robustness} summarizes the two frozen policies and their arithmetic mean. Without perturbations, mean success was $96.65\%$ on this test set. With position noise, success was $97.60\%$ at $\sigma_p=0.5$ mm, but fell to $53.10\%$ at $1$ mm and $0.05\%$ at $2$ mm. Thus, noise on the scale of the position-success tolerance substantially degraded reaching performance. Under the tested nonzero orientation-noise levels ($0.0025$--$0.020$ rad per component), mean success ranged from $97.30\%$ to $99.15\%$. Small nonmonotonic increases do not establish a general benefit from measurement noise.

Force-only mean success remained $96.65\%$ at 0, 5, and 10 N and $96.60\%$ at 20 N. This characterizes the policy together with its underlying velocity servo; it does not isolate disturbance rejection by the learned policy. With simultaneous 20-N force and position noise, mean success decreased from $96.25\%$ at $0.75$ mm to $88.60\%$, $80.60\%$, $67.05\%$, and $51.70\%$ at $0.85$, $0.90$, $0.95$, and $1.00$ mm, respectively. At $1.00$ mm, the individual policies achieved $35.8\%$ and $67.6\%$, exposing sensitivity to the training seed. Because the single- and combined-perturbation sweeps used different task blocks and the combined sweep held force magnitude fixed, these results do not identify a noise--force interaction. They also do not establish a robustness advantage over other methods.
\par\endgroup

\begin{figure*}[t]
\centering
\includegraphics[width=\textwidth]{figures/fig_robustness.pdf}
\caption{\textcolor{blue}{Frozen Full policies at $1$ mm/$0.025$ rad under (a) position measurement noise, (b) orientation measurement noise, (c) randomly redirected external tool force, and (d) position noise with a 20-N tool force. Colored curves show two training seeds; the dashed curve is their mean. Each policy uses 1000 episodes per condition. Noise magnitudes are standard deviations per translation or rotation-vector component. Panels (a) and (b) share the zero-noise evaluation. Tasks are paired within each sweep; panel (d) uses a different task block. Success uses the true simulated state.}}
\label{fig:robustness}
\end{figure*}
